\section{Integral Dependence and Valuations}

\begin{exercise}
	$f^*(V(I))=\left\{f^{-1}(\mathfrak{q}):\mathfrak{q}\supseteq I\right\}$. Use \Cref{5.10} to proof $f^*(V(I))\supseteq V\left(f^{-1}(I)\right)$.
\end{exercise}

\begin{exercise}
	
\end{exercise}

\begin{exercise}
	Use the operator $-\otimes_A c^n$ to get the result.
\end{exercise}

\begin{exercise}
	If $\left(\frac{1}{x+1}\right)^n+\sum\frac{f_i\left(x^2-1\right)}{g_i\left(x^2-1\right)(x+1)^i}=0$, multiply both side by $(x+1)^n$ and take $x=-1$ to get contradiction. 
\end{exercise}

\begin{exercise}
	(i) Element-wise. (ii) Use \Cref{5.10} to proof that the contraction of maximal ideals in $B$ are maximal in $A$
\end{exercise}

\begin{exercise}
	Let $b_i\in B_i$, $f_i\in A[x]$ monic, $f_i(b_i)=0$. Let $f=f_1\cdots f_n$, then $f(b_1,\cdots,b_n)=0$.
\end{exercise}

\begin{exercise}
	Let $x\in B$ and $x$ integral over $A$. Let $x^n+a_{n-1}x^{n-1}+\cdots+a_0=0$ is the equation of the least degree. Analyze $x\left(x^{n-1}+a_{n-2}x^{n-2}+\cdots+a_1\right)=-a_0$.
\end{exercise}

\begin{exercise}
	(ii) Proof that there is a ring $D\supseteq B$ such that $f$ factors into linear factors in $D$. Use Kronecker's construction.
\end{exercise}

\begin{exercise}
	Hint.
\end{exercise}

\begin{exercise}\label{ex5.10}
	(a) $\Rightarrow$ (b): Let $f^*(V(\mathfrak{q}_1))=V(I)$, so $I\subseteq\mathfrak{p}_1\subseteq\mathfrak{p}_2$, and we have $\mathfrak{p}_2\in V(I)=f^*(V(\mathfrak{q}_1))$. (a') $\Rightarrow$ (b'): Hint.
\end{exercise}

\begin{exercise}\label{ex5.11}
	Use \Cref{ex3.18,ex5.10}.
\end{exercise}

\begin{exercise}\label{ex5.12}
	Note that $G$ has identity element so $x$ is a root of the polynomial $\prod_{\sigma\in G}\left(t-\sigma(x)\right)$. The coefficients are symmetric polynomials of $\sigma_i$. We don't know how $G$ acts on $\left(S^{-1}A\right)^G$, so we define $\sigma(x/s)=\sigma(x)/\sigma(s)$.
\end{exercise}

\begin{exercise}
	
\end{exercise}

\begin{exercise}
	$\sigma\in G$ fixes $A$.
\end{exercise}

\begin{exercise}
	
\end{exercise}

\begin{exercise}
	
\end{exercise}

\begin{exercise}
	Hint. Using Zariski's lemma is the easiest way.
\end{exercise}

\begin{exercise}[\textcolor{red}{Zariski's lemma}]
	
\end{exercise}

\begin{exercise}
	Omitted.
\end{exercise}

\begin{exercise}
	Hint.
\end{exercise}

\begin{exercise}
	Hint.
\end{exercise}

\begin{exercise}\label{ex5.22}
	Hint.
\end{exercise}

\begin{exercise}\label{ex5.23}
	(i) $\Rightarrow$ (ii): Note that $f:A\rightarrow A/\operatorname{ker}f$. Use \Cref{1.1}. (ii) $\Rightarrow$ (iii): Consider $A/\mathfrak{p}$. (iii) $\Rightarrow$ (i): If there is a prime ideal $\mathfrak{p}$ which is not an intersection of maximal ideals, the Jacobson radical $\mathfrak{J}$ of $A/\mathfrak{p}$ is not zero. Pick $f\in\mathfrak{J}$, then $(A/\mathfrak{p})_f\neq0$. The contraction of a maximal ideal of $(A/\mathfrak{p})_f$ is maximal among all prime ideals which doesn't contain $f$, but by (iii), the intersection has to contain $f$.
\end{exercise}

\begin{exercise}
	(i) Use \Cref{ex5.23}(iii) and going-up theorem. (ii) Consider $B/\mathfrak{q}$ and $A/\mathfrak{q}^c$. Then use \Cref{ex5.22}, \Cref{ex5.23}(i) and \Cref{1.1}.
\end{exercise}

\begin{exercise}
	Hint.
\end{exercise}

\begin{exercise}
	(2) $\Rightarrow$ (3): $U_1\cap X_0=U_2\cap X_0$ $\Rightarrow$ $U_1^c\cap X_0=U_2^c\cap X_0$, then take closure. (i) $\Rightarrow$ (ii) and (ii) $\Rightarrow$ (iii): Consider $V(\mathfrak{a})\cap X_f\cap\operatorname{Max}(A)$. If all maximal ideals in $V(\mathfrak{a})$ contains $f$, $V(\mathfrak{a})\cap X_f=\varnothing$ by \Cref{ex5.23}(iii). (iii) $\Rightarrow$ (i): If not Jacobson, there is a prime ideal $\mathfrak{p}$ such that $\mathfrak{p}\subsetneq\bigcap_{\mathfrak{q}\supseteq\mathfrak{p}}\mathfrak{q}$. Let $f\in\left.\bigcap_{\mathfrak{q}\supseteq\mathfrak{p}}\mathfrak{q}\middle\backslash\mathfrak{p}\right.$, then $X_f\cap V(\mathfrak{p})=\{\mathfrak{p}\}$.
\end{exercise}

\begin{exercise}
	Use \Cref{5.21} to proof maximal elements are valuation rings. Conversely, if $(A,\mathfrak{m})\subsetneq(A^\prime,\mathfrak{m}^\prime)$, pick $x\in A^\prime\backslash A$. Then $x^{-1}\in\mathfrak{m}$ because it is not a unit in $A$.
\end{exercise}

\begin{exercise}
	(1) $\Rightarrow$ (2): Pick $a\in\mathfrak{a}\backslash\mathfrak{b}$, and any $b\in\mathfrak{b}$. Consider $a/b$ and $b/a$. (2) $\Rightarrow$ (1): Consider $(a)\subseteq(b)$ or $(b)\subseteq(a)$, we have $a=xb$ or $b=ya$, so $y=x^{-1}$.
\end{exercise}

\begin{exercise}
	
\end{exercise}


